\documentclass{article}
\usepackage{amsmath,amssymb}
\usepackage{xurl}
\usepackage[colorlinks=true,linkcolor=blue,citecolor=blue,urlcolor=blue]{hyperref}
% Shared commands for notes in docs/paper-gaps/.

\DeclareUnicodeCharacter{2080}{\ifmmode{}_{0}\else\textsubscript{0}\fi}
\DeclareUnicodeCharacter{2081}{\ifmmode{}_{1}\else\textsubscript{1}\fi}
\DeclareUnicodeCharacter{2082}{\ifmmode{}_{2}\else\textsubscript{2}\fi}
\DeclareUnicodeCharacter{2099}{\ifmmode{}_{n}\else\textsubscript{n}\fi}

\newcommand{\Lean}{\textsc{Lean}}
\ExplSyntaxOn
\NewDocumentCommand{\leanid}{m}{
  \ifmmode
    \textsf{\detokenize{#1}}
  \else
    \group_begin:
    \urlstyle{sf}
    \tl_set:Nn \l_tmpa_tl {#1}
    \tl_replace_all:Nnn \l_tmpa_tl {\_} {_}
    \exp_args:NV \path \l_tmpa_tl
    \group_end:
  \fi
}
\ExplSyntaxOff
\providecommand{\tr}{\operatorname{tr}}
\newcommand{\tnleanIssueBase}{https://github.com/LionSR/TNLean/issues/}
\newcommand{\tnleanPrBase}{https://github.com/LionSR/TNLean/pull/}
\newcommand{\tnleanDocBase}{https://sirui-lu.com/TNLean/docs/}
\newcommand{\ghissue}[1]{\href{\tnleanIssueBase#1}{GitHub issue \##1}}
\newcommand{\ghpr}[1]{\href{\tnleanPrBase#1}{PR \##1}}
\newcommand{\doclink}[1]{\href{\tnleanDocBase#1.html}{\path{#1}}}

% Dirac notation and the identity operator, as in blueprint/src/macros/common.tex.
% \providecommand keeps note-local definitions authoritative.
\usepackage{dsfont}
\providecommand{\ket}[1]{|#1\rangle}
\providecommand{\bra}[1]{\langle #1|}
\providecommand{\braket}[2]{\langle #1 | #2 \rangle}
\providecommand{\ketbra}[2]{|#1\rangle\!\langle #2|}
\providecommand{\Id}{\mathds{1}}
\providecommand{\one}{\mathds{1}}
\providecommand{\C}{\mathbb{C}}
\providecommand{\MN}[1]{M_{#1}(\C)}
\providecommand{\MD}{M_D(\C)}

% Verdict marker: \gapnote{<kind>}{<status>}. Kind is the policy
% classification (clarification, local-correction, scope-restriction,
% unfaithful, false-source, open-gap); status is open, wip (elimination
% underway), resolved, or historical. Machine-read by the site index;
% prints a verdict line.
\newcommand{\gapnote}[2]{\par\noindent\textbf{Verdict:} #1 (#2).\par}

\providecommand{\leanid}[1]{\textsf{\detokenize{#1}}}
\providecommand{\gapnote}[2]{\par\noindent\textbf{Verdict:} #1 (#2).\par}
\title{The Error of the Localized Filtered Parts in Hastings' Lemma~1}
\author{}
\date{October 22, 2026}
\begin{document}
\maketitle
\gapnote{local-correction}{resolved}

\section{The source statement}

In~\cite{Hastings2007AreaLaw} (arXiv:0705.2024v4, \path{law16.tex},
lines~85--143), the chain has sites $1,\dots,N$, the Hamiltonian is
$H=\sum_iH_{i,i+1}$ with $H_{i,i+1}$ supported on the sites $i,i+1$ and
$\Vert H_{i,i+1}\Vert\le J$, $\Delta E$ is the spectral gap above the unique
ground state $\Psi_0$, $v$ and $\xi_C$ are the velocity and length scale of the
Lieb--Robinson bound, and $\xi=\max(2v/\Delta E,\xi_C)$. For integers $a\le b$,
$X_{a,b}$ is the set of sites $i$ with $a\le i\le b$. Lemma~1 concerns a site
$j$ and a length $l$.

In the appendix of the source (lines~774--792), the Hamiltonian is split as
$H=H_L+H_B+H_R$, where $H_L$ is the sum of the terms $H_{i,i+1}$ with $i\le
j-l/3$, $H_R$ the sum of those with $i\ge j+1+l/3$, and $H_B$ the rest, and constants are added to the three parts so that their
ground-state expectations vanish (lines~781--782). Each
part is filtered with a Gaussian in time,
$\tilde H_L^0=(\Delta E/\sqrt{2\pi q})\int H_L(t)e^{-(t\Delta E)^2/2q}\,dt$
with $q=(l/3)\Delta E/(2v)$, and the source states (lines~816--828) that, by the
Lieb--Robinson bound, the filtered parts can be approximated by operators
$M_L,M_B,M_R$ with $M_L-H_L$ supported on $X_{j-2l/3,j}$ and
\begin{align}
    \Vert M_L-H_L\Vert\le\mathcal O(\Delta E^{-1}J^2e^{-l/3\xi}),
    \label{eq:hastings07_localized_parts_source}
\end{align}
and similarly for $M_B$ and $M_R$. From eq.~\texttt{tighter} (line~797), which bounds
$|\tilde H_L^0\Psi_0|$, it then concludes
\begin{align}
    |M_L\Psi_0|\le\mathcal O(\Delta E^{-1}J^2e^{-l/3\xi}).
    \label{eq:hastings07_localized_parts_ground_state}
\end{align}

\section{The point at issue}

The printed bound on $\Vert M_L-H_L\Vert$ cannot hold in general: $\tilde
H_L^0-H_L$ is the Gaussian average of $H_L(t)-H_L$, which is not small, and
$M_L-H_L$ is supported near the cut while $\tilde H_L^0-H_L$ need not be small
there. The next step, $|M_L\Psi_0|\le|\tilde H_L^0\Psi_0|+\Vert M_L-\tilde
H_L^0\Vert$, uses the distance to the filtered part. The intended statement is
\begin{align}
    \Vert M_L-\tilde H_L^0\Vert\le\mathcal O(\Delta E^{-1}J^2e^{-l/3\xi}),
    \label{eq:hastings07_localized_parts_intended}
\end{align}
with $M_L-H_L$ supported near the cut.

\section{The statement proved}

For any operator $W$, write $\tau_t(W)=W(t)=e^{iHt}We^{-iHt}$, let $\tilde W$ be
the Gaussian average of $\tau_t(W)$ with the weight used for $\tilde H_L^0$, and
set $M=W+E_K(\tilde W-W)$, where $E_K$ is the normalized partial trace onto a
set of sites $K$. Then $M-W$ acts on $K$, and if the evolved
commutator $\tau_s(i[H,W])$ is within $\varepsilon$ of its partial trace onto
$K$ for $|s|\le T$,
\begin{align}
    \Vert\tilde W-M\Vert\le T\varepsilon
      +2\Vert[H,W]\Vert\sqrt{2h/\pi}\,e^{-T^2/(2h)},\qquad h=q/\Delta E^2.
    \label{eq:hastings07_localized_parts_proved}
\end{align}
The bound involves $[H,W]$ and not $W$: by the Duhamel formula,
$\Vert\tau_t(W)-W\Vert\le|t|\,\Vert[H,W]\Vert$. For $W=H_L$ the commutator is
supported near the cut, and the chain Lieb--Robinson bound supplies an
$\varepsilon$ independent of the length of the chain.\footnote{%
\leanid{TNLean.Lattice.norm_gaussianFilteredTerm_sub_localized_le} in
\path{TNLean/Lattice/GaussianFilteredIncrement.lean}.}

\section{Conclusion}

The correction replaces the printed bound (\ref{eq:hastings07_localized_parts_source})
by the intended one (\ref{eq:hastings07_localized_parts_intended}), which follows from
(\ref{eq:hastings07_localized_parts_proved}) with $W=H_L$. The bound used by the rest
of the proof, (\ref{eq:hastings07_localized_parts_ground_state}), follows as in
the source.

\bibliographystyle{alpha}
\bibliography{references}
\end{document}
